\section{Introduction}
The field of video generation is advancing at a remarkable pace, catalyzed by the advent of diffusion models. Seminal works such as Sora~\citep{openai2024sora}, Wan~\citep{wan2025}, Hunyuan-DiT~\citep{kong2024hunyuanvideo}, and Veo~\citep{deepmind_veo} are progressively closing the gap between generated content and reality. Despite this progress, a formidable challenge remains: the majority of state-of-the-art models are confined to generating short-form videos, typically capped at 5-10 seconds. This constraint is inherent to the architectural design of the underlying Diffusion Transformers (DiT)~\citep{peebles2023scalable}, the inherently non-streaming and non-causal nature of the vanilla DiT architecture poses a significant challenge to achieving temporal scalability,

A promising avenue for transcending this limitation lies in shifting from bidirectional diffusion architectures to autoregressive, streaming-based models. One such approach, Diffusion Forcing~\citep{chen2024diffusion, kim2024fifo}, applies heterogeneous noise schedules across frames to enable sequential generation. However, the combinatorial complexity of noise scheduling often leads to training instability and has proven difficult to scale~\cite{chen2025skyreels,he2025matrix}.
A more tractable strategy involves predicting the next frame or chunk from a clean context, with KV caching emerging as a key mechanism for enabling performant, real-time streaming. For instance, CausVid~\citep{yin2025causvid} proposes a method to distill a bidirectional teacher model into a streaming student model using heterogeneous distillation. However, its reliance on overlapping frames for temporal consistency and a pronounced train-inference mismatch often results in over-exposure artifacts. The Self-Forcing~\citep{huang2025self} method mitigates the over-exposure issue by aligning the training and inference distributions. While this sets a new benchmark for short-form video quality, its capacity remains bottlenecked by the fixed-duration teacher model. Consequently, when tasked with generating content beyond this intrinsic temporal window (e.g., $>$10 seconds), the model's visual quality degrades precipitously. 


A primary challenge limiting the quality of autoregressive long-video generation models is a significant \textit{training-inference misalignment}. This misalignment manifests in two principal ways. First, a \textit{temporal mismatch} occurs: during training, models generate short clips of up to 5 seconds—the maximum horizon of the teacher model—whereas at inference, they must generate videos of significantly greater length. Second, error accumulation caused by \textit{supervision misalignment} during long-horizon generation. In training, the teacher model provides abundant supervision for every frame within the short clip. This intensive guidance, however, means the student model is rarely exposed to the compounding errors that naturally arise in long rollouts, leaving it ill-equipped to handle them. As a result, generation quality rapidly deteriorates beyond the 5-second training horizon, often collapsing into static or stalled content.

In this paper, we introduce \textbf{Self-Forcing++}, which directly targets the above two issues. Building upon the observation from previous works \citep{bruce2024genie,alonso2024diffusion,li2025hunyuan} that a teacher model, despite its own 5-second generation limit, possesses rich knowledge for correcting errors in quality-degraded videos due to its training on a vast video corpus. We leverage this insight by extending the student's generation horizon far beyond 5 seconds (up to 100 seconds in our experiments). This process intentionally produces candidate long videos that contain accumulated errors. To enable the student model to handle these errors, we then re-inject noise into these degraded rollouts and apply distribution-matching distillation with the strong teacher model, a process combined with a long-horizon rolling KV cache and windowed sampling. This strategy teaches the student to recover from degraded states and sustain high-quality, coherent video generation over extended durations.
Experimental results demonstrate that our method can scale video generation up to 100 seconds, a 20$\times$ increase over the baseline, while maintaining high visual quality. By scaling up computation through extended training, our method is capable of generating videos up to 4 minutes and 15 seconds\footnote{The maximum number of latent frames Wan2.1-T2V-1.3B supports is 1024, since we generate videos in a trunk size of 3, the maximum length we can reach is 1023 which is 99.9\% of maximally supported length 1024.}, utilizing 99.9\% of the base model's positional embedding capacity and representing a 50$\times$ improvement over the baseline.
Furthermore, our investigation revealed that the widely used VBench~\citep{huang2024vbench} benchmark exhibits a bias that favors over-exposed and degraded frames when evaluating long videos, undermining the reliability of its results. To remedy this, we propose a new metric, Visual Stability, designed to systematically capture both quality degradation and over-exposure in long video generation. Our work paves the way for building more robust and reliable long video generation models.
Our contributions are summarized as follows:

\begin{itemize}[leftmargin=*]
\item \underline{\textit{Identifying Horizon Scaling Bottlenecks}}: We reveal the primary obstacle to extending the generation horizon of autoregressive models: a dual mismatch in temporality and supervision during training versus inference. This insight provides a clear target for overcoming previous limitations on the generation length.

\item \underline{\textit{A Simple Solution}}: We propose a simple training framework, named Self-Forcing++. By generating beyond the teacher's horizon and correcting the student model on its own long, error-accumulated rollout trajectories, Self-Forcing++ extends high-quality video generation to 100 seconds, far surpassing previous state-of-the-art methods without reusing overlapping frames.

\item \underline{\textit{SOTA Performance and Horizon Scalability}}: Self-Forcing++ achieves state-of-the-art (SOTA) performance in long-video generation across a range of durations (e.g., 10s, 50s, 100s). Furthermore, we discover a significant scaling property: by scaling the training computation, our model's generation capability extends to multiple minutes, a feat previously considered out of reach.
\end{itemize}